\section{The day-weighted clock}
\label{sec:clkclock}

We now know the market keeps its own clock, and that a desk can only
carry a forecast of its schedule.  This section builds that forecast
in the simplest currency a trading desk speaks: \emph{days}.

Each scheduled event --- an earnings report, a central-bank decision,
a court ruling --- is recorded as a pair $(t_e, N_e)$: its date
$t_e$, a calendar year fraction like any expiry's, and its size
$N_e$, the number of \emph{extra} equivalent days the event is worth.
Mind the word extra: an earnings day with $N_e=4$ counts as five
ordinary days of variance --- its own one, plus four extra.  The
walk's event day in \cref{sec:clkwalk}, five variance units in one
calendar day, is in this currency an event of size $N_e=4$; the
five-unit-wide strip of \cref{fig:clkwalk}(b) and the four-day jump
about to appear in \cref{fig:clkclock}(a) describe the same event in
the two currencies.  Given a calendar of such events, the variance
time of maturity $t$ is defined by counting days:
\begin{equation}
\boxed{\;\tau_{\rm days}(t)\;=\;365\,t
\;+\!\!\sum_{t_e\le t}\!N_e,
\qquad
\tau(t)\;=\;\frac{\tau_{\rm days}(t)}{365}.\;}
\label{eq:clkclock}
\end{equation}
The sum runs over events at or before $t$: an expiry pays for exactly
the events it contains.  With an empty calendar, $\tau(t)=t$ --- the
convention every previous chapter ran under, now visible as the
special case.

Compute one by hand.  Take a 14-day expiry and a single event worth
$N_e=\MacClkClockExtraDays$ extra days, scheduled
\MacClkClockEventDay\ days out.  The event is inside the expiry, so
\[
\tau_{\rm days}=14+\MacClkClockExtraDays=18,
\qquad
\tau=\frac{18}{365}=0.0493,
\]
against a calendar $t=14/365=0.0384$.  Eighteen variance days will
elapse in fourteen calendar days.  For an expiry of 9 days --- one
that \emph{dodges} the event --- the sum is empty and
$\tau_{\rm days}=9$: nothing changes.  That is the whole mechanism;
everything else in the chapter is consequences.

\Cref{fig:clkclock}(a) draws \eqref{eq:clkclock} for this calendar.
The blue line is $\tau_{\rm days}$ against calendar days: it tracks
the dashed calendar diagonal exactly until day \MacClkClockEventDay,
jumps by \MacClkClockExtraDays\ days there, and runs \emph{parallel}
to the diagonal afterwards --- an event shifts the clock once and is
never paid for again.  (The dotted variant is the normalized clock,
explained at the end of this section; at this zoom it hugs the plain
one.)  Panel~(b) is read after the proposition below.

What may a change of clock do, and what may it never do?  The answer
deserves to be a proposition, because it is the safety case for
everything that follows.  Recall two terms from Chapter~2, one line
each: a smile is \emph{butterfly-valid} when its implied density is
nonnegative --- the Durrleman check $g_D\ge0$ of \eqref{eq:durrleman}
--- and two expiries are in \emph{calendar order} when total variance
does not decrease with maturity at fixed log-moneyness
(\cref{sec:calendar}).

\begin{proposition}[Relabeling time: invariance and covariance]
\label{prop:clkgauge}
Let $\tau(\cdot)$ be any strictly increasing relabeling of maturity
with $\tau(0)=0$ --- in particular the clock \eqref{eq:clkclock}.
Then:
\begin{enumerate}[label=(\roman*)]
\item option prices, and with them every total variance $w(k)$ and
the terminal law of each expiry, are unchanged;
\item butterfly validity of each smile and calendar order across
smiles are unchanged;
\item the volatility reading is \emph{covariant} --- it does not stay
fixed but transforms with the clock, by a known factor:
$\sigma=\sqrt{w/\tau}=\sigma_{\rm cal}\sqrt{t/\tau}$;
\item the accrual rate between consecutive expiries is covariant in
the same sense, by the interval's clock ratio:
$\Delta w/\Delta\tau
=(\Delta w/\Delta t)\cdot(\Delta t/\Delta\tau)$.
\end{enumerate}
\end{proposition}

\begin{proof}
A relabeling attaches new time coordinates to the same expiries; no
random variable, no expectation, and no price changes, which is (i).
Butterfly validity is a property of one expiry's law, unchanged by
(i); calendar order compares the same pairs $w(k,T_1)\le w(k,T_2)$ in
the same order because the relabeling is strictly increasing, which
is (ii).  For (iii) and (iv), divide the same invariant quantities
($w$, $\Delta w$) by the relabeled denominators and cancel:
$\sqrt{w/\tau}=\sqrt{w/t}\cdot\sqrt{t/\tau}$, and likewise for the
ratios of increments.
\end{proof}

Everything a desk can trade sits in the invariant column; everything
the clock touches is a \emph{reading}.  Part~(iv) looks like bare
cancellation, and that is exactly its content: the increments
$\Delta w$ carry all the information, and a clock change re-divides
them without adding any --- a fact \cref{sec:clkinverse} will lean
on.  In particular, installing an event can never move a price and
never create an arbitrage --- not because an implementation is
careful, but because the clock enters the mathematics in exactly one
place: as a denominator.

Part (iii) is worth turning around, because it is the formula desks
feel.  Fix an expiry $t$ and insert one event of $N_e$ extra days
before it.  Total variance $w$ is pinned by the price, so the clock
reading drops by exactly the ratio of the two clocks:
\begin{equation}
\sigma
=\sigma_{\rm cal}\sqrt{\frac{t}{\tau}}
=\sigma_{\rm cal}\sqrt{\frac{365\,t}{365\,t+N_e}} .
\label{eq:clkcrush}
\end{equation}
Keep the direction straight, once, on the hand example.  The clock
reading is the \emph{smaller} number: \eqref{eq:clkcrush} multiplies
the calendar reading by $\sqrt{14/18}$, taking $34.0\%$ down to
$30.0\%$.  Turned over, the calendar reading is the clock reading
\emph{lifted} by
$\sqrt{18/14}=\MacClkClockRatioTwoWeeks$: the same pair of numbers,
the same price.  How big is the effect in general?  For an event small
against the maturity ($N_e\ll365\,t$), expand the square root with
the binomial approximation $\sqrt{1+x}\approx1+\tfrac{x}{2}$, whose
error is of order $x^2$:
\[
\frac{\sigma_{\rm cal}}{\sigma}
=\sqrt{1+\frac{N_e}{365\,t}}
\;\approx\;1+\frac{N_e}{2\cdot365\,t}.
\]
The lift of the calendar reading scales like $N_e$ over \emph{twice
the calendar days to expiry}: four extra days lift a one-year reading
by about half a percent of itself ($4/730$), and a two-week reading
by roughly a seventh ($4/28$; the exact lift is
$\MacClkClockRatioTwoWeeks-1$, a touch less).  Events are a
short-maturity phenomenon --- exactly where desks watch implied
volatility most closely.

Now read \cref{fig:clkclock}(b), which plots the lift direction: the
blue curve is $\sigma_{\rm cal}/\sigma=\sqrt{\tau/t}$ across
expiries, for the same one-event calendar.  Expiries before day
\MacClkClockEventDay\ dodge the event: the ratio is exactly one.  The
expiry \emph{on} the event date pays the full toll,
$\sqrt{14/10}=\MacClkClockPeakRatio$ --- the annotated peak.  Past it
the ratio decays toward one as the event becomes a smaller and
smaller fraction of the maturity.  The dashed curve is the
small-event approximation just derived: generous at the peak, where
$N_e$ is anything but small against $365\,t$, still visibly above
the exact curve out to about a month, and indistinguishable beyond.

\begin{figure}[htbp]
\centering
\paperfig{\textwidth}{fig_clk_clock.pdf}
\caption{The day-weighted clock, one event of
\MacClkClockExtraDays\ extra days at day \MacClkClockEventDay.
(a)~$\tau_{\rm days}(t)$: the calendar diagonal, a jump of
\MacClkClockExtraDays\ days at the event, parallel growth after; the
normalized variant (dotted) rescales the whole line by
\MacClkClockNormFactor.  (b)~The reading ratio
$\sigma_{\rm cal}/\sigma=\sqrt{\tau/t}$ across expiries: one before
the event, a peak of \MacClkClockPeakRatio\ at the event-day expiry,
decay like $1+N_e/(2\cdot365\,t)$ (dashed) beyond it.}
\label{fig:clkclock}
\end{figure}

One counting choice completes the clock.  As defined,
\eqref{eq:clkclock} makes a year with events \emph{longer} than 365
variance days: the events add variance on top of an unchanged
baseline.  Some desks prefer the opposite convention: the year's
total budget is what it is, and events say \emph{when} it accrues,
not that there is more of it.  The normalized clock implements that
by rescaling all days --- ordinary and event alike --- by one factor
that restores the annual budget:
\begin{equation}
\tau_{\rm days}(t)\;\longmapsto\;
\tau_{\rm days}(t)\cdot
\frac{365}{365+\sum_{t_e\le1}N_e},
\label{eq:clknorm}
\end{equation}
where the sum in the denominator collects the events of the first
year ($t_e\le1$).  For the worked calendar the factor is
$365/369=\MacClkClockNormFactor$ --- the dotted line of
panel~(a).  Under \eqref{eq:clknorm} the one-year variance time, and
with it the one-year volatility reading, exactly match a no-event
year.  Sub-year readings shift slightly: the rescaling makes every
day worth a little less than before, so a maturity that dodges all
events now counts marginally fewer variance days, while the events'
\emph{relative} weight within the year is exactly what it was.
Neither convention is more correct:
un-normalized treats events as extra variance against a
one-per-day baseline, normalized treats the annual budget as known
and the events as its scheduled lumps.  What matters --- the book's
recurring refrain --- is that the choice is stated.  The chapter's
contract table records it, and the whole rest of the chapter runs
un-normalized.

We now own the clock: a desk-sized object --- a list of dated day
counts --- that turns the erratic readings of \cref{sec:clkboard}
into steady ones without touching a single price.  The next section
watches it work through the week every option trader knows best.
